\documentclass[11pt]{article}
\usepackage[T1]{fontenc}
\usepackage{lmodern}
\usepackage[margin=1in,top=.84in,bottom=.88in]{geometry}
\usepackage{amsmath,amssymb,amsthm,booktabs,microtype}
\usepackage[hidelinks]{hyperref}
\newtheorem{theorem}{Theorem}
\newtheorem{lemma}[theorem]{Lemma}
\newtheorem{corollary}[theorem]{Corollary}
\newcommand{\ph}{\varphi}
\newcommand{\fl}[1]{\lfloor #1\rfloor}
\newcommand{\fr}[1]{\{#1\}}
\title{Prime columns and primitive rows\\in the Wythoff array}
\author{Blake Stone mathematics repository\\[-2pt]\small Proof developed with OpenAI Codex}
\date{}
\begin{document}
\maketitle
\begin{abstract}
Every column $k$ of the Wythoff array contains a proportion
$\ph^{-(k+1)}$ of the primes, where $\ph$ is the golden ratio and $k$ starts at $1$.
Yet one explicit primitive row consists entirely of composite numbers.
This proves Kimberling's column conjecture and refutes his row conjecture.
We identify the counterexample as a tail of Vsemirnov's composite Fibonacci-like
sequence and give a congruence class of initial pairs containing a positive
density of primitive, prime-free rows. We also prove that precisely
$6/\pi^2$ of rows, in natural density, are primitive, and that their
longest consecutive run has length five.
\end{abstract}

\section{Two prime conjectures}
Write $\ph=(1+\sqrt5)/2$, $\beta=\ph-1=1/\ph$,
$F_0=0$, $F_1=1$ and $F_{k+2}=F_{k+1}+F_k$. Morrison's Wythoff array
\cite{morrison} has entries
\begin{equation}\label{W}
 W(m,k)=F_{k+1}\fl{m\ph}+(m-1)F_k,\qquad m,k\ge1 .
\end{equation}
We also use \eqref{W} at $k=0$ as an auxiliary initial value.
A row is primitive when its consecutive entries are coprime.
Kimberling asked whether every column contains infinitely many primes, and
whether every primitive row does \cite{primes}.

\begin{theorem}\label{primes}
For each $k\ge1$,
\[
 \lim_{X\to\infty}
 \frac{\#\{\ell\le X:\ell\text{ prime and }\ell=W(m,k)\text{ for some }m\ge1\}}{\pi(X)}
 =\beta^{k+1}.
\]
The row $M=5\,381\,684\,723\,742\,272\,612\,570$ is primitive and contains
no primes.
\end{theorem}

\subsection*{The interval for a column}
For a positive integer $n$ set $f_n=\fr{n\beta}$. We claim that
\begin{equation}\label{window}
 n\text{ belongs to column }k
 \quad\Longleftrightarrow\quad
 f_n\in
 \begin{cases}
 (\beta^{k+2},\beta^k) & k\text{ even},\\
 (1-\beta^k,1-\beta^{k+2}) & k\text{ odd}.
 \end{cases}
\end{equation}
To see this, take $b=m-1$, $a=\fl{m\ph}$, and write
$n=bF_k+aF_{k+1}$, $h=bF_{k-1}+aF_k$. From
$F_{t+1}\beta-F_t=(-\beta)^t\beta$ we obtain
\[
 n\beta-h=(-\beta)^k x,\qquad x=\beta a-b=1-\beta f_m .
\]
Since $m>0$, we have $\beta^2<x<1$, giving the forward implication.
Conversely, either open interval in \eqref{window} gives an integer $h$
and an $x$ in this range with $n\beta=h+(-\beta)^k x$. The vectors
$(F_k,F_{k-1})$, $(F_{k+1},F_k)$ are a basis of $\mathbb Z^2$, so
$(n,h)$ has integral coefficients $b,a$ in this basis. The identity above
forces $x=\beta a-b$. If $b\le-1$, the inequality $x<1$ forces $a\le0$
and hence $n\le0$. Thus $m=b+1\ge1$. Multiplication of
$\beta^2<x<1$ by $\ph$ gives $\beta<a-\ph b<\ph$,
which implies $a<m\ph<a+1$. This proves \eqref{window}.

Both intervals have length $\beta^k(1-\beta^2)=\beta^{k+1}$.
Vinogradov's theorem gives uniform distribution of $\fr{\ell\beta}$
as $\ell$ runs through the primes, because $\beta$ is irrational
\cite[\S5]{vinogradov}. Applying it to \eqref{window} proves the asserted
relative density and infinitude in each column.

\subsection*{A prime-free primitive row}
Vsemirnov \cite{vsemirnov} proved that the Fibonacci-like sequence with seeds
$A_0=106276436867$, $A_1=35256392432$ has no prime terms.
Morrison's array lists the positive Fibonacci sequences up to a shift
\cite{morrison,orderings}. To locate this example explicitly, recurrence gives
\begin{align*}
 A_{53}&=M-1=5381684723742272612569,&
 A_{54}&=8707748799751085288495,\\
 A_{55}&=14089433523493357901064,&
 A_{56}&=22797182323244443189559.
\end{align*}
Set $a_0=A_{54}$ and $u_0=A_{53}$. The exact integer inequalities
\[
 (2a_0-M)^2\le 5M^2<(2a_0+2-M)^2,\qquad 2a_0-M\ge0
\]
give $\fl{M\ph}=a_0$. Therefore $W(M,1)=A_{55}$, $W(M,2)=A_{56}$,
and the shared recurrence gives $W(M,k)=A_{k+54}$ for all $k\ge1$.
The Euclidean algorithm gives $\gcd(u_0,a_0)=1$, and the gcd of successive
terms is invariant under the recurrence. This completes Theorem~\ref{primes}.

For an independently checkable certificate put $T_k(a,u)=F_{k+1}a+F_k u$.
The following are Vsemirnov's covering data with each index residue shifted
by $-54$. Each triple $(s,t,r)$ certifies that $s$ divides the term at every
column $k\equiv r\pmod t$.
\[
\begin{array}{rrr@{\qquad}rrr@{\qquad}rrr}
\toprule
s&t&r&s&t&r&s&t&r\\\midrule
3&4&1&47&16&3&1103&48&27\\
2&3&1&19&18&14&181&90&26\\
5&5&0&61&15&3&41&20&4\\
7&8&7&23&24&11&541&90&8\\
17&9&2&107&36&26&2521&60&54\\
11&10&2&31&30&6&&&\\\bottomrule
\end{array}
\]
Indeed $s\mid F_t$ and $s\mid T_r(a_0,u_0)$ for each triple. The identity
$T_{k+t}=F_{t-1}T_k+F_tT_{k+1}$ propagates the divisibility.
The periods divide $720$, and the listed classes cover every residue
$0,\ldots,719$. Since $T_k(a_0,u_0)\ge a_0>2521$, all divisors are proper.
The Lean proof verifies the table and coverage by integer arithmetic.

\section{Consecutive primitive rows}
Let $u=m-1$, $a=\fl{(u+1)\ph}$ and $G(u)=\gcd(u,a)$.
The auxiliary term $W(m,0)=a$ and the next term $a+u$ have gcd $G(u)$.
Consecutive pairs throughout the row have this same gcd. Entries A332937
and A332938 ask for the largest run with $G=1$ and its limiting proportion
\cite{primitive}.

\begin{theorem}
The largest number of consecutive primitive rows is five.
Furthermore the row of index $\ell+1$ is primitive for every prime $\ell>3$.
\end{theorem}
\begin{proof}
Among six consecutive row indices, take the three odd ones $m,m+2,m+4$.
For such indices $u$ is even. At least one of
$\fl{m\ph},\fl{(m+2)\ph},\fl{(m+4)\ph}$ is even. Indeed, both consecutive
increments are at least three, and the total increment is at most seven.
If all three numbers were odd, each increment would be at least four.
This contradiction gives a nonprimitive row in any run of six.
For rows $36$ through $40$ the five floor values are $58,59,61,63,64$.
Their gcds with $35,36,37,38,39$, respectively, are all one.

For a prime $\ell>3$ we have $\ell\ge5$ and
$\ell<\fl{(\ell+1)\ph}<2\ell$. Therefore $\ell$ cannot divide that floor.
This proves the observation of Karttunen in A332937.
\end{proof}

\section{The distribution of row gcds}
\begin{theorem}\label{density}
For every integer $g\ge1$ the row indices with gcd $g$ have natural density
\[
 \frac{6}{\pi^2 g^2}.
\]
In particular the primitive-row density is $6/\pi^2=0.60792710\ldots$.
\end{theorem}
This also follows from Estermann's theorem on primitive points in strips
of irrational slope \cite{estermann}. Indeed, put $u=gv$ and $a=gw$.
The points $(v,w)$ lie in a strip of vertical width $1/g$ with $v\le N/g$.
We give a self-contained counting reduction for this particular slope. Its
uniform bound will also handle residue restrictions on both coordinates.

Let $C_d(N)=\#\{1\le u\le N:d\mid G(u)\}$. For $d\ge2$, write $u=dv$.
There is an integer $w$ with $dw=\fl{(dv+1)\ph}$ precisely when
\begin{equation}\label{strip}
 \beta<d(w-v\ph)<\ph.
\end{equation}
The endpoints are strict because $(dv+1)\ph$ is irrational.
As $0<\beta/d<\ph/d<1$, such a $w$ must be $\fl{v\ph}+1$. Thus
\[
 1-\ph/d < \fr{v\beta} < 1-\beta/d .
\]
The interval has length $1/d$, and $1\le v\le\fl{N/d}$. Equidistribution
gives $C_d(N)/N\to1/d^2$, also true for $d=1$.

\begin{lemma}\label{majorant}
For $d,N\ge1$ we have $C_d(N)\le10N/d^2$.
\end{lemma}
\begin{proof}
If $v\ge1$, $w\in\mathbb Z$ and $|w-v\ph|\le2$, the nonzero integer
$(w-v\ph)(w+v\beta)=w^2-vw-v^2$ has absolute value at least one. Also
$|w+v\beta|\le2+v(\ph+\beta)\le5v$. Thus $1\le5v|w-v\ph|$.
For a point of \eqref{strip}, $|w-v\ph|<2/d$ gives $d^2\le10u$.
For two such points at $u_1<u_2$, subtract \eqref{strip} to obtain
$|(w_2-w_1)-(v_2-v_1)\ph|<1/d$. The same norm argument gives
$d^2\le5(u_2-u_1)$. These inequalities also apply at $d=1$, since
\eqref{strip} itself still holds for points counted. Hence
$u\mapsto\fl{10u/d^2}$ injects the counted points into the integers
from $1$ to $\fl{10N/d^2}$.
\end{proof}

The divisor identity for the Möbius function gives
$\mathbf1_{G(u)=g}=\sum_{gd\mid G(u)}\mu(d)$ for $u>0$.
The lemma bounds each normalized summand in absolute value by
$10/(gd)^2$. Dominated convergence now yields
\[
 \lim_{N\to\infty}\frac1N\#\{1\le u\le N:G(u)=g\}
 =\sum_{d\ge1}\frac{\mu(d)}{(gd)^2}
 =\frac{1}{g^2\zeta(2)}=\frac{6}{\pi^2g^2}.
\]
The omitted first row has no effect on density.

\section{A positive-density family of prime-free primitive rows}
Let $P=1821895895860356790898731230$ be the product of the $17$ displayed
primes. Consider the rows with $u\equiv u_0\pmod P$ and
$a=\fl{(u+1)\ph}\equiv a_0\pmod P$. Each $T_k(a,u)$ is congruent to
$T_k(a_0,u_0)$ modulo $P$. For $a>2521$ the certificate therefore makes
every entry composite.

\begin{corollary}\label{family}
The primitive rows in these two residue classes, with $a>2521$, have
natural density
\[
 \frac1{P^2}\prod_{\ell\ {\rm prime},\ \ell\nmid P}(1-\ell^{-2})>0 .
\]
\end{corollary}
\begin{proof}
For any modulus $L\ge1$ and prescribed residues $r,s\bmod L$, the conditions
$u\equiv r$, $a\equiv s\pmod L$ have joint natural density $1/L^2$.
Indeed, for $u=r+Lv$, the second condition says that
$\{v\ph+(r+1)\ph/L\}$ belongs to $[s/L,(s+1)/L)$, taking
$0\le s<L$. This rotated irrational orbit is equidistributed. Among the
$1/L$ fraction of indices in the first progression it selects a fraction
$1/L$.

If $(d,P)=1$, the Chinese remainder theorem combines the two prescribed
residues with $d\mid u,a$ into one pair modulo $Pd$. Its density is
$1/(Pd)^2$. If $(d,P)>1$, the conditions are incompatible since
$\gcd(u_0,a_0)=1$. For the restricted count of primitive rows apply
Möbius inversion again. The absolute contribution of each $d$, after
division by $N$, is bounded by $10/d^2$ by Lemma~\ref{majorant}.
Interchanging the limit and the sum gives
$P^{-2}\sum_{(d,P)=1}\mu(d)/d^2$ and hence the stated product.
It is positive by convergence of $\sum_{\ell}\ell^{-2}$. Discarding the
finitely many small values $a\le2521$ leaves this limit unchanged.
\end{proof}
This density concerns the specified class of rows. It supplies a lower bound
on the lower density of all prime-free primitive rows.

\paragraph{Formal scope.}
Lean proves the column intervals, the transfer of prime equidistribution to
column densities and infinitude, the explicit primitive composite row, the
congruence certificate, the run bound and witness, the prime-offset observation
and Theorem~\ref{density}. The transfer takes Vinogradov's published theorem
as one explicit hypothesis. The limit calculation for the restricted family
in Corollary~\ref{family} is given above.

\begin{thebibliography}{9}
\bibitem{morrison} D. R. Morrison, A Stolarsky array of Wythoff pairs,
\emph{A Collection of Manuscripts Related to the Fibonacci Sequence},
Fibonacci Association, 1980, 134--136.
\bibitem{orderings} C. Kimberling, Orderings of the set of all positive Fibonacci
sequences, \emph{Applications of Fibonacci Numbers} 5 (1993), 405--416.
\href{https://doi.org/10.1007/978-94-011-2058-6_39}{DOI}.
\bibitem{primes} C. Kimberling, OEIS \href{https://oeis.org/A360377}{A360377}
and \href{https://oeis.org/A360378}{A360378} (2023). See also the conjectural
row assumption in \href{https://oeis.org/A333086}{A333086} (2020).
\bibitem{primitive} C. Kimberling, OEIS \href{https://oeis.org/A332937}{A332937}
and \href{https://oeis.org/A332938}{A332938} (2020), with A. Karttunen's
prime-offset observation in A332937 (2025).
\bibitem{vinogradov} J.-M. De Koninck and I. K\'atai, Spacings between selected
prime divisors of an integer, \emph{Bull. Aust. Math. Soc.} 107 (2023), 403--415.
\href{https://doi.org/10.1017/S0004972722000880}{DOI}.
\bibitem{vsemirnov} M. Vsemirnov, A new Fibonacci-like sequence of composite
numbers, \emph{J. Integer Seq.} 7 (2004), Article 04.3.7.
\href{https://cs.uwaterloo.ca/journals/JIS/VOL7/Vsemirnov/vsem5.pdf}{Article}.
\bibitem{estermann} T. Estermann, On the number of primitive lattice points in a
parallelogram, \emph{Canad. J. Math.} 5 (1953), 456--459.
\href{https://doi.org/10.4153/CJM-1953-051-6}{DOI}.
\end{thebibliography}
\end{document}
